Calcul de l'auteur. Base : 777 629 ventes comportant au moins un logement recensées dans DVF en 2024 (hors Alsace-Moselle et Mayotte). Élasticités : -7 \% de transactions par point de droit (Fritzsche et Vandrei) ; -15 \% pour un droit de 1,1 \% (Dachis et al.) ; -37 \% de déménagements pour 2 points (Hilber et Lyytikäinen, effet sur la mobilité et non sur les ventes) ; hausse de 0,7 point en 2014 en France : ventes +26 \% le mois précédent, -14 \% par mois les trois mois suivants, puis aucun effet (Bérard et Trannoy). Droits sur l'ancien : 5,8 points avant la hausse de 2025.\par Lecture : à gauche, effet d'une baisse d'un point, proche des variations observées dans les études ; la seule estimation française ne trouve pas d'effet durable. À droite, les barres hachurées prolongent ces élasticités jusqu'à la suppression des droits, très au-delà de ce qui a été observé ; le modèle de Bontemps et al. simule le passage de 14 \% à 6 \% de coûts de transaction (droits, frais et commissions) compensé par un impôt sur la propriété.
